\documentclass[11pt]{article}
\usepackage[margin=1in]{geometry}
\usepackage{amsmath,amssymb,amsthm,mathtools}
\usepackage{graphicx}
\usepackage[hidelinks]{hyperref}
\usepackage{microtype}

\newtheorem{theorem}{Theorem}
\newtheorem{lemma}{Lemma}
\newtheorem{remark}{Remark}
\newcommand{\R}{\mathbb R}
\newcommand{\tr}{\operatorname{tr}}
\newcommand{\Span}{\operatorname{span}}
\newcommand{\ip}[2]{\langle #1,#2\rangle}

\title{One size-two atom suffices for the quartic tracial moment problem on $Y^2=1$}
\author{Candidate full solution packet for arXiv:2001.11614\\
Agent \texttt{agent\_lane\_03} \quad Model GPT5.6}
\date{9 August 2026}

\begin{document}
\maketitle

\noindent\textbf{Status.} \texttt{full\_solution\_likely\_valid}; subject to expert review and a broader novelty search.

\section{Source problem}

Bhardwaj and Zalar study normalized bivariate quartic tracial moment
sequences and representing measures by pairs of real symmetric matrices,
using normalized matrix trace \cite{BZ}.  Their Conjecture~1, on PDF page 12,
asks whether, under the column relation $Y^2=1$, the following are
equivalent: existence of an arbitrary representing measure; existence of a
measure using only scalar and size-two atoms; and existence of a measure of
type $(m,1)$, i.e. with exactly one size-two atom.

\begin{center}
\includegraphics[width=\textwidth]{figures/open_problem_crop.png}
\end{center}

The preceding Theorem~4.1 in \cite{BZ} proves this when the original measure
has no atoms larger than size three.  The theorem below removes that
restriction and proves the full conjecture.

\section{Definitions and a scalar moment fact}

A finite representing measure can be encoded as a finite-dimensional real
$*$-algebra
\[
  \mathcal A=\bigoplus_{j=1}^N M_{n_j}(\R)
\]
with the tracial state
\[
  \tau(T)=\sum_{j=1}^N \rho_j\,\frac{1}{n_j}\tr(T_j),
  \qquad \rho_j>0,\quad \sum_j\rho_j=1.
\]
The given moments are $\beta_w=\tau(w(X,Y))$.  Conversely, such a direct
sum is precisely a finite tracial representing measure.

We use the following standard degree-four Hamburger fact (equivalently, the
one-variable flat moment theorem; see, for example, \cite{Akh}).

\begin{lemma}[Scalar flat completion]\label{lem:scalar}
Let $s_0,s_1,s_2,s_3\in\R$ and put
\[
 H=\begin{pmatrix}s_0&s_1\\s_1&s_2\end{pmatrix},
 \qquad g=\binom{s_2}{s_3}.
\]
If $H\succ0$, then setting $s_4=g^TH^{-1}g$ produces a positive
semidefinite flat Hankel matrix
\[
 \begin{pmatrix}H&g\\g^T&s_4\end{pmatrix},
\]
and $(s_0,\ldots,s_4)$ has a representing measure on $\R$ with at most two
atoms.  The same conclusion holds with $H^\dagger$ when $H\succeq0$ and
$g\in\operatorname{ran}H$; in rank one the representing measure is a point
mass.
\end{lemma}

\section{Proof intuition}

Diagonalizing the involution $Y$ splits $X$ into two diagonal corners and
one off-diagonal block:
\[
 X=\begin{pmatrix}A&B\\B^*&C\end{pmatrix},
 \qquad Y=\begin{pmatrix}1&0\\0&-1\end{pmatrix}.
\]
All failure of commutativity is measured, through degree four, by $B$.
We choose a single size-two atom whose off-diagonal entry absorbs the total
energy $\tau(BB^*)$, while its diagonal entries are the $BB^*$- and
$B^*B$-weighted means of $A$ and $C$.  This choice cancels every mixed term
through degree three.  What remains on the two signs of $Y$ is ordinary
one-variable moment data.

The apparent obstruction is the fourth moment: making the size-two atom too
light forces its off-diagonal entry to be large.  The exact minimum fourth
moments of the two scalar residuals are Schur complements.  Orthogonal
projection in the two corner $L^2$ spaces turns their sum into a one-variable
function $G$.  Two applications of Cauchy--Schwarz and one elementary
two-variable inequality show $G\geq0$ at an explicit parameter.  Since
$G\to-\infty$ at the other end, continuity gives a parameter with $G=0$;
both scalar completions are then flat.  This simultaneously handles every
atom in the original measure and leaves exactly one size-two atom.

\section{The theorem}

\begin{theorem}\label{thm:main}
Let $\beta=(\beta_w)_{|w|\leq4}$ be a normalized bivariate real tracial
moment sequence whose moment matrix satisfies the column relation $Y^2=1$.
If $\beta$ has a representing measure, then it has a representing measure of
type $(m,1)$.  If $\beta$ is noncommutative, one may take $m\leq4$.
Consequently all three conditions in Conjecture~1 of \cite{BZ} are
equivalent.
\end{theorem}

\begin{proof}
Encode an arbitrary finite representing measure by $(\mathcal A,\tau,X,Y)$
as above.  The column relation is inherited atomwise: indeed,
\[
 0=\tau\bigl((Y^2-1)^2\bigr)
   =\sum_j\rho_j\frac1{n_j}\tr\bigl((Y_j^2-I_{n_j})^2\bigr),
\]
and every summand is nonnegative because $Y_j$ is symmetric.  Thus
$Y_j^2=I_{n_j}$ for every atom, and $Y^2=1$ in $\mathcal A$.  Put
\[
 e_+=\frac{1+Y}{2},\qquad e_-=\frac{1-Y}{2},
\]
and, relative to these two corners, write
\[
 X=\begin{pmatrix}A&B\\B^*&C\end{pmatrix}.
\]
All traces below are taken with $\tau$; block formulas are interpreted
componentwise in the direct sum.

\subsection*{Step 1: the moment identities}

Set
\[
 D=BB^*,\qquad D'=B^*B,\qquad
 w=\tau(D)=\tau(D'),\qquad
 S=\tau(D^2)=\tau((D')^2).
\]
For $0\leq k\leq3$ let
\[
 p_k=\frac{\beta_{X^k}+\beta_{X^kY}}2=\tau(e_+X^k),\qquad
 q_k=\frac{\beta_{X^k}-\beta_{X^kY}}2=\tau(e_-X^k).
\]
Direct block multiplication gives
\begin{align*}
 p_0&=\tau(e_+),&q_0&=\tau(e_-),\\
 p_1&=\tau(A),&q_1&=\tau(C),\\
 p_2&=\tau(A^2)+w,&q_2&=\tau(C^2)+w,\\
 p_3&=\tau(A^3)+2\tau(AD)+\tau(CD'),
 &q_3&=\tau(C^3)+\tau(AD)+2\tau(CD').
\end{align*}
Moreover,
\begin{align}
 \beta_{X^4}
 &=\tau(A^4)+\tau(C^4)
   +4\{\tau(A^2D)+\tau(C^2D')+\tau(ABCB^*)\}+2S,\label{eq:x4}\\
 \Delta:=\beta_{X^2}-\beta_{XYXY}&=4w.\label{eq:delta}
\end{align}

If $w=0$, then $B=0$ in the tracial $L^2$ norm and $X$ commutes with $Y$.
Spectral measures of $A$ and $C$ give a scalar representing measure.
Replacing part of any scalar atom by the identical scalar pair on a
two-dimensional space gives a representation of type $(m,1)$ if that exact
format is desired.  Hence assume $w>0$, so also $S>0$.

\subsection*{Step 2: the unique natural size-two atom}

Define the weighted means
\[
 a=\frac{\tau(AD)}w,\qquad c=\frac{\tau(CD')}w.
\]
For a parameter $\lambda>0$, let
\[
 \widetilde X=\begin{pmatrix}a&b\\b&c\end{pmatrix},\qquad
 \widetilde Y=\begin{pmatrix}1&0\\0&-1\end{pmatrix},\qquad
 b^2=\frac w\lambda,
\]
and give this size-two atom density $2\lambda$.  Since normalized trace on
$M_2$ is one half of ordinary trace, the atom contributes $\lambda$ times
the ordinary two-dimensional trace.  In particular its contribution to
$\Delta$ is $4\lambda b^2=4w$, exactly \eqref{eq:delta}.

The plus-corner contributions through degree three are
\[
 \lambda,\quad \lambda a,\quad \lambda(a^2+b^2),\quad
 \lambda\{a^3+b^2(2a+c)\},
\]
and the minus-corner formulas interchange $a$ and $c$.  By the definitions
of $a,c$ and $\lambda b^2=w$, subtraction therefore leaves
\begin{equation}\label{eq:residual-cubic}
 r_k^+=\tau(A^k)-\lambda a^k,qquad
 r_k^-=\tau(C^k)-\lambda c^k,qquad 0\leq k\leq3.
\end{equation}
Thus the remaining problem is scalar on each sign of $Y$.

\subsection*{Step 3: positivity and exact scalar fourth moments}

Work first in the plus corner, with inner product
$\ip{R}{T}=\tau(RT)$.  Let $L_A=\Span\{e_+,A\}$, let $Q_A$ be the
orthogonal projection of $A^2$ onto $L_A$, and put
\[
 R_A=A^2-Q_A,\qquad E_A=\|R_A\|_2^2.
\]
Evaluation of an affine-linear polynomial at $a$ is represented on $L_A$ by
the projection of $D/w$.  Write
\[
 \frac Dw=P_{L_A}\!\left(\frac Dw\right)+z_A,qquad
 \kappa_A=\left\|P_{L_A}\!\left(\frac Dw\right)\right\|_2^2,qquad
 s_A^2=\|z_A\|_2^2=\frac S{w^2}-\kappa_A.
\]
If $Q_A=\alpha e_++\gamma A$, define
\[
 h_A=\alpha+\gamma a,qquad \delta_A=h_A-a^2,qquad
 V_A=\tau((A-ae_+)^2D).
\]
The same definitions, with $C,D',e_-$, give
$R_C,E_C,z_C,\kappa_C,s_C,h_C,\delta_C,V_C$.

Cauchy--Schwarz shows that for every affine-linear $\ell$,
\[
 |\ell(a)|^2
 =\left|\tau\left(\ell(A)\frac Dw\right)\right|^2
 \leq \tau(\ell(A)^2)\frac S{w^2}.
\]
Consequently, whenever
\begin{equation}\label{eq:lambda-range}
 \frac1\lambda\geq\frac S{w^2},
\end{equation}
both order-one Hankel matrices associated to \eqref{eq:residual-cubic} are
positive semidefinite.  Also $w^2\leq S\tau(e_+)$ and
$w^2\leq S\tau(e_-)$, so \eqref{eq:lambda-range} implies
$\lambda\leq\min\{\tau(e_+),\tau(e_-)\}$ and $2\lambda\leq1$.

Assume first that the base plus Hankel matrix is nonsingular.  A
Sherman--Morrison calculation says that the least fourth moment completing
$(r_0^+,\ldots,r_3^+)$ is
\begin{equation}\label{eq:qA}
 q_A(\lambda)=\tau(Q_A^2)-\lambda a^4
 +\frac{\lambda\delta_A^2}{1-\lambda\kappa_A}.
\end{equation}
Indeed, the base order-one Hankel matrix is the Gram matrix of
$e_+,A$; its cubic cross-vector represents $Q_A$, and subtracting
$\lambda(1,a)^T(1,a)$ gives \eqref{eq:qA}.  The identical formula holds for
$q_C(\lambda)$.  If a base Hankel matrix has rank one, the corresponding
corner operator is scalar in $L^2(\tau)$ and the same formula holds by
continuity, with the zero quotient interpreted as zero.  At a singular
endpoint the necessary range condition is likewise the limiting one.

Two identities will control the apparent poles.  Orthogonality of $R_A$ to
$L_A$ gives
\begin{equation}\label{eq:delta-identity}
 \delta_A=\frac{V_A}{w}-\ip{R_A}{z_A},
 \qquad
 \delta_C=\frac{V_C}{w}-\ip{R_C}{z_C}.
\end{equation}
For example,
$\tau(R_AD/w)=\tau(A^2D/w)-h_A=V_A/w-\delta_A$, and the left side is
$\ip{R_A}{z_A}$.

\subsection*{Step 4: the fourth-moment surplus}

Put
\[
 T=V_A+V_C+\tau((A-ae_+)B(C-ce_-)B^*).
\]
The ordinary trace of $\widetilde X^4$ is
\[
 a^4+c^4+4b^2(a^2+ac+c^2)+2b^4.
\]
Using \eqref{eq:x4}, \eqref{eq:qA}, and $\lambda b^2=w$, the amount of
fourth moment left after paying for the size-two atom and the two minimal
scalar completions is exactly
\begin{align}
 G(\lambda)
 &=E_A+E_C+4T+2S-\frac{2w^2}{\lambda}\notag\\
 &\quad-\frac{\lambda\delta_A^2}{1-\lambda\kappa_A}
       -\frac{\lambda\delta_C^2}{1-\lambda\kappa_C}.
 \label{eq:G-lambda}
\end{align}

Write
\[
 \frac1\lambda=\frac S{w^2}+v,\qquad v\geq0.
\]
Then \eqref{eq:G-lambda} becomes
\begin{equation}\label{eq:G-v}
 G(v)=E_A+E_C+4T-2w^2v
 -\frac{\delta_A^2}{s_A^2+v}
 -\frac{\delta_C^2}{s_C^2+v}.
\end{equation}
By \eqref{eq:delta-identity} and Cauchy--Schwarz in the direct sums
$L^2\oplus\R$,
\begin{align}
 \frac{\delta_A^2}{s_A^2+v}
 &\leq E_A+\frac{V_A^2}{w^2v},&
 \frac{\delta_C^2}{s_C^2+v}
 &\leq E_C+\frac{V_C^2}{w^2v}.
 \label{eq:CS}
\end{align}
Here $v>0$; the zero case is read as a limit.

Let $U=(A-ae_+)B$ and $V=B(C-ce_-)$.  Then
$V_A=\|U\|_2^2$, $V_C=\|V\|_2^2$, and the last term in $T$ is
$\ip{U}{V}$.  Hence
\[
 T\geq V_A+V_C-\sqrt{V_AV_C}.
\]
For all $x,y\geq0$,
\begin{equation}\label{eq:elementary}
 4(x+y-\sqrt{xy})\geq2\sqrt2\sqrt{x^2+y^2};
\end{equation}
after squaring, the difference is $2(\sqrt x-\sqrt y)^4$.

If $V_A^2+V_C^2>0$, choose
\[
 v_0=\frac{\sqrt{V_A^2+V_C^2}}{\sqrt2,w^2}.
\]
Equations \eqref{eq:G-v}--\eqref{eq:elementary} give
\[
 G(v_0)\geq
 4T-\left(2w^2v_0+\frac{V_A^2+V_C^2}{w^2v_0}\right)
 =4T-2\sqrt2\sqrt{V_A^2+V_C^2}\geq0.
\]
If $V_A=V_C=0$, use $v_0=0$.  Then
\eqref{eq:delta-identity} gives
$\delta_A^2/s_A^2\leq E_A$ and
$\delta_C^2/s_C^2\leq E_C$ (with the usual zero convention), so again
$G(v_0)\geq0$.

On the other hand $G(v)\to-\infty$ as $v\to\infty$.  It follows by
continuity that there is $\widehat v\geq v_0$ such that
$G(\widehat v)=0$.  Choose
\[
 \lambda=\left(\frac S{w^2}+\widehat v\right)^{-1}.
\]
For this choice, the remaining fourth moment is exactly
$q_A(\lambda)+q_C(\lambda)$.

\subsection*{Step 5: construct the scalar atoms}

Give the plus residual in \eqref{eq:residual-cubic} fourth moment
$q_A(\lambda)$ and the minus residual fourth moment $q_C(\lambda)$.
Lemma~\ref{lem:scalar} produces a scalar measure on the line $y=1$ and one
on the line $y=-1$, each with at most two atoms.  Together with the single
size-two atom $(\widetilde X,\widetilde Y)$ of density $2\lambda$, these
measures match $p_k,q_k$ for $0\leq k\leq3$, match $\beta_{X^4}$, and match
$\Delta$.  Cyclicity and $Y^2=1$ reduce every word of degree at most four to
these data, so the resulting measure represents all of $\beta$.

Thus a noncommutative $\beta$ has a measure of type $(m,1)$ with $m\leq4$.
The implications from type $(m,1)$ to a measure using atoms of sizes at most
two and then to an arbitrary representing measure are immediate.  This proves
all equivalences in the source conjecture.
\end{proof}

\section{Verification and limitations}

The accompanying script \texttt{code/verify\_random\_reduction.py} checks the
block identities, residual cubic matching, exact surplus formula, and its
proved lower bound on 5,000 seeded random finite tracial direct sums.  This is
a regression test only; the proof above is analytic.

The delicate review points are: the coefficient $4$ in \eqref{eq:x4}; the
Sherman--Morrison term in \eqref{eq:qA}; and the limiting convention when a
corner Hankel matrix has rank one.  In the latter case the corner operator is
scalar in $L^2(\tau)$, so its residual is a point mass and the formulas are
literal limits of the nonsingular case.

A bounded novelty check was performed on 9 August 2026.  Exact local-corpus
and arXiv/web searches used arXiv:2001.11614, the paper title and authors, and
the phrases ``$Y^2=1$'', ``type $(m,1)$'', and ``atoms of size at most 2''.
They found the source and earlier adjacent work, but no later resolution.
This does not establish novelty; a broader expert literature search remains
necessary.

\section{Human-review recommendation}

High priority.  The argument is short once expressed in a finite tracial
direct sum, it answers the exact conjecture, and it strengthens the conclusion
to at most four scalar atoms in the noncommutative case.  A reviewer should
recompute \eqref{eq:x4} and \eqref{eq:G-lambda}, then check the two singular
endpoint cases against Lemma~\ref{lem:scalar}.

\begin{thebibliography}{9}
\bibitem{BZ}
A. Bhardwaj and A. Zalar,
\emph{The tracial moment problem on quadratic varieties},
J. Math. Anal. Appl. 498 (2021), 124936;
arXiv:\href{https://arxiv.org/abs/2001.11614}{2001.11614}.

\bibitem{Akh}
N. I. Akhiezer,
\emph{The Classical Moment Problem and Some Related Questions in Analysis},
Oliver \& Boyd, 1965.
\end{thebibliography}

\end{document}
